\subsection{Vision Transformer}
\label{sec:vision_transformer}

Arsitektur ekstraksi representasi visual penelitian ini mengadopsi ViT-Base (vit-base-patch16-224) sebagaimana diilustrasikan pada Figure~\ref{fig:3}. Citra wajah masukan berukuran 224 × 224 piksel dengan 3 saluran warna dibagi menjadi 196 patch spasial non-overlapping $\mathbf{x}_p^i$ berukuran 16 × 16 piksel. Setiap patch diproyeksikan secara linier melalui matriks $\mathbf{E}$ ke ruang laten berdimensi $D = 768$, lalu dirangkaikan dengan token kelas $\mathbf{x}_{\text{class}}$ dan embedding posisi $\mathbf{E}_{\text{pos}}$ sesuai \eqref{eq:1}, di mana $\mathbf{z}_0$ merupakan vektor sekuens token awal masukan encoder:
\begin{equation}
\mathbf{z}_0 = [\mathbf{x}_{\text{class}}; \, \mathbf{x}_p^1\mathbf{E}; \, \dots; \, \mathbf{x}_p^{196}\mathbf{E}] + \mathbf{E}_{\text{pos}}
\label{eq:1}
\end{equation}

\begin{figure}[htbp]
\centering
\includegraphics[width=\columnwidth]{vit.png}
\caption{Architecture of the ViT Backbone and Patch Projection.}
\label{fig:3}
\end{figure}

Vektor sekuens token awal $\mathbf{z}_0$ dialirkan ke tumpukan 12 layer transformer encoder identik. Setiap layer encoder $\ell$ memproses masukan $\mathbf{z}_{\ell-1}$ menjadi representasi antara $\mathbf{z}'_\ell$ melalui mekanisme MHSA setelah operasi Layer Normalization (LN) sesuai \eqref{eq:2}, di mana $\text{LN}(\cdot)$ menyatakan normalisasi lapisan dan $\mathbf{z}_{\ell-1}$ adalah representasi luaran dari layer sebelumnya:
\begin{equation}
\mathbf{z}'_\ell = \text{MHSA}(\text{LN}(\mathbf{z}_{\ell-1})) + \mathbf{z}_{\ell-1}
\label{eq:2}
\end{equation}
Proses tersebut dilanjutkan menuju representasi luaran $\mathbf{z}_\ell$ melalui blok Multi-Layer Perceptron (MLP) dua lapis beraktivasi GeLU berdasarkan \eqref{eq:3}, di mana $\text{MLP}(\cdot)$ merupakan jaringan perceptron multilapis dan $\mathbf{z}'_\ell$ adalah representasi antara berresidual:
\begin{equation}
\mathbf{z}_\ell = \text{MLP}(\text{LN}(\mathbf{z}'_\ell)) + \mathbf{z}'_\ell
\label{eq:3}
\end{equation}
Mekanisme self-attention dengan 12 attention heads memetakan relasi spasial global antarpatch wajah secara langsung tanpa pembatasan receptive field lokal.

Untuk mempertahankan efisiensi komputasi dan mencegah variabilitas pelatihan ulang, ketiga model backbone ViT dibekukan (frozen) sebagai penyedia task-associated representations secara offline. Model ViT-Face (skutaada/VIT-VGGFace) menangkap representasi terkait geometri biometrik wajah, ViT-Emotion (dima806/facial\_emotions\_image\_detection) menghasilkan representasi terkait ekspresi wajah, dan ViT-Age (dima806/facial\_age\_image\_detection) mengekstraksi representasi terkait estimasi usia wajah. Dari setiap model domain, vektor fitur $\mathbf{f}_{\text{domain}} \in \mathbb{R}^{768}$ diekstraksi dari representasi token $[\text{CLS}]$ pada layer encoder terakhir $L$ ($\mathbf{z}_L^0$) setelah operasi LN sesuai \eqref{eq:4}, di mana $\mathbf{z}_L^0$ merupakan vektor representasi token pertama pada layer ke-$L$:
\begin{equation}
\mathbf{f}_{\text{domain}} = \text{LN}(\mathbf{z}_L^0) \in \mathbb{R}^{768}
\label{eq:4}
\end{equation}

Ketiga representasi laten domain tunggal tersebut digabungkan menjadi vektor fusi tri-domain $\mathbf{z}_{\text{tri}} \in \mathbb{R}^{2304}$ melalui operasi konkatenasi fitur ($\oplus$) dari vektor $\mathbf{f}_{\text{face}}$, $\mathbf{f}_{\text{emotion}}$, dan $\mathbf{f}_{\text{age}}$ sebagaimana dirumuskan pada \eqref{eq:5}, di mana $\mathbf{z}_{\text{tri}}$ menyatakan vektor representasi terfusi tri-domain:
\begin{equation}
\mathbf{z}_{\text{tri}} = \mathbf{f}_{\text{face}} \oplus \mathbf{f}_{\text{emotion}} \oplus \mathbf{f}_{\text{age}} \in \mathbb{R}^{2304}
\label{eq:5}
\end{equation}

Eksplorasi sistematis mencakup tujuh skema ablasi fitur, yang meliputi tiga konfigurasi domain tunggal berdimensi 768, tiga konfigurasi domain ganda berdimensi 1.536, serta satu konfigurasi tri-domain berdimensi 2.304 seperti dirangkum pada Table~\ref{tab:II}.

\begin{table}[htbp]
\caption{Multi-Domain Feature Fusion and Ablation Configurations.}
\label{tab:II}
\centering
\begin{tabular}{clcc}
\toprule
\# & Configuration & Domain Category & Dimension \\
\midrule
1 & Face & Single-Domain & 768 \\
2 & Emotion & Single-Domain & 768 \\
3 & Age & Single-Domain & 768 \\
4 & Emotion $\oplus$ Face & Dual-Domain & 1,536 \\
5 & Face $\oplus$ Age & Dual-Domain & 1,536 \\
6 & Emotion $\oplus$ Age & Dual-Domain & 1,536 \\
7 & Face $\oplus$ Emotion $\oplus$ Age & \textbf{Tri-Domain (Proposed)} & \textbf{2,304} \\
\bottomrule
\end{tabular}
\end{table}
